\documentclass[10pt,a4paper]{amsart}
\usepackage[margin=19mm]{geometry}
\usepackage{amsmath,amssymb,mathtools}
\usepackage[scr=rsfs,cal=euler]{mathalfa}
\usepackage{xfrac}
\usepackage[unicode,hidelinks,bookmarks=false]{hyperref}
\newcommand{\ie}{i.e.\ }
\newcommand{\cf}{cf.\ }
\newcommand{\resp}{respectively\ }
\newcommand{\Subsection}{\subsection}
\providecommand{\nobreakdash}{\nobreak\hbox{-}}
\setlength{\emergencystretch}{2em}
\multlinegap0pt
\usepackage{mathtools}
\usepackage{amssymb}
\usepackage[shortlabels]{enumitem}
\usepackage{xcolor}
\usepackage[normalem]{ulem}
\newtheorem{proposition}{Proposition}
\title{Sharp bounds and geometric properties of the first non trivial Steklov Neumann Eigenvalue}
\author{Sagar Basak, Gloria Paoli, Rossano Sannipoli, Sheela Verma}
\date{Source: arXiv:2603.25448v1 (CC BY 4.0)}
\begin{document}
\maketitle

\noindent\emph{Source and licence.} Sharp bounds and geometric properties of the first non trivial Steklov Neumann Eigenvalue. Version arXiv:2603.25448v1 (CC BY 4.0), distributed by arXiv under the Creative Commons Attribution 4.0 International licence, http://creativecommons.org/licenses/by/4.0/, as recorded in the arXiv licence metadata for this exact version and shown on the abstract page https://arxiv.org/abs/2603.25448v1. Full licence notice: \texttt{licenses/arxiv-2603.25448-CC-BY-4.0.txt}.\par\smallskip

\noindent\textit{Editorial scope.} Counted unit: the article's proposition proposition (source0.tex lines 1048--1050). The article's complete original proof of it is delivered with it (source0.tex lines 1051--1104, one \texttt{proof} environment of 54 lines). No mathematics is written, repaired or re-derived: the statement and the proof are the article's own lines, in the article's own notation, and no retained line is substituted or re-flowed.  No further source-local context is needed: every cross-reference of every retained line resolves inside the two delivered spans.  Nothing was omitted from any retained span, so the package carries no omission record.  No retained line contains a citation, so the wrapper carries no bibliography and no bibliography entry.  No retained line contains a figure environment, an image-inclusion command or drawing code, so no image is delivered and the package declares no asset dependency.  Editorial formatting only, in addition: an AMS \texttt{amsart} preamble that restates the article's own declarations and macros used by the retained lines, namely the environments \texttt{proposition} on separate counters, together with the standard packages those lines need.  The retained environments print in this document as Proposition 1 in order of appearance, whereas the article prints the same items with the numbering of its own full document.  The complete native source of the article as distributed in the arXiv e-print for this exact version is bundled unchanged as \texttt{sources/197212/source0.tex}.
\begin{proposition}
    There exists a sequence $\{\Omega_\varepsilon\}_{\varepsilon}\subseteq \mathbb R^2$ of open, bounded sets with Lipschitz boundary, such that $\mu_1(\Omega_\varepsilon)\to 0$, as $\varepsilon\to 0^+$.
\end{proposition}

\begin{proof}
Let $\varepsilon>0$ be small enough. Let us consider  the open  rectangle centered at the origin of sides $\varepsilon$ and $\varepsilon^3$ 
\begin{equation*}
C_{\varepsilon}= \bigg(-\frac{\varepsilon}{2},\frac{\varepsilon}{2}\bigg)\times \bigg(-\frac{\varepsilon^3}{2},\frac{\varepsilon^3}{2}\bigg)
\end{equation*}
 and let us define  the two balls
\begin{equation*}
    B_{i,\varepsilon}=B_{r_i(\varepsilon)}(x_i(\varepsilon)),  \qquad i=1,2,
\end{equation*}
with radii $r_1(\varepsilon)=r_2(\varepsilon)= \sqrt{(1+\varepsilon/2)^2+\varepsilon^6/4}$, centered at the points $x_1(\varepsilon)=(-1-\varepsilon,0)$ and $x_2(\varepsilon)= (1+\varepsilon,0)$, respectively.
Let us consider the following sequence of perforated dumbells $\{\Omega_\varepsilon\}\subset \mathbb{R}^2 $, defined as follows
\begin{equation*}
	\Omega_\varepsilon = B_{1,\varepsilon}\cup C_{\varepsilon}\cup (B_{2,\varepsilon}\setminus \overline{B_{R_1}(x_2(\varepsilon))}),
\end{equation*}
where $B_{R_1}(x_2(\varepsilon))$ is the concentric ball contained in $B_{r_2(\varepsilon)}(x_2(\varepsilon))$ of radius $0<R_1<1$. Let us consider the following continuous function
\begin{equation*}
	v(x,y)= \begin{cases}
	\sin \bigg(\displaystyle\frac{2\pi x}{\varepsilon}\bigg) & \text{in} \, C_{\varepsilon}\\
	0 & \text{elsewhere}.
	\end{cases}
\end{equation*}
Let us stress that $v(x,y)$ is a good test function for the first non zero Steklov--Neumann eigenvalue, since
\begin{equation*}
    \int_{\partial \Omega_\varepsilon\setminus \partial B_{R_1}(x_2(\varepsilon))}v\,d\mathcal{H}^{n-1}  = \int_{\partial C_\varepsilon}v\,d\mathcal{H}^{n-1} = 2\int_{-\frac{\varepsilon}{2}}^{\frac{\varepsilon}{2}}\sin\bigg(\displaystyle\frac{2\pi x}{\varepsilon}\bigg)\,dx = -\frac{\varepsilon}{\pi}\cos \bigg(\displaystyle\frac{2\pi x}{\varepsilon}\bigg)\Bigg|_{-\frac{\varepsilon}{2}}^{\frac{\varepsilon}{2}}=0.
\end{equation*}
The denominator in the Rayleigh quotient of $\mu_1(\Omega_\varepsilon)$ becomes

\begin{equation*}
\begin{split}
\int_{\partial \Omega_\varepsilon}v^2 \,d\mathcal{H}^{n-1}&= \int_{\partial C_\varepsilon}v^2\,d\mathcal{H}^{n-1} = 2\int_{-\frac{\varepsilon}{2}}^{\frac{\varepsilon}{2}}\sin^2 \bigg(\displaystyle\frac{2\pi x}{\varepsilon}\bigg)\,dx\\
&=4\int_0^{\frac{\varepsilon}{2}} \sin^2\bigg(\displaystyle\frac{4\pi x}{\varepsilon}\bigg)\,dx
=4\displaystyle\int_0^{\frac{\varepsilon}{2}}\frac{1-\cos\bigg(\displaystyle\frac{4\pi x}{\varepsilon}\bigg)}{2}\,dx=\varepsilon.
\end{split}
\end{equation*}
Moreover, since
\begin{equation*}
	|\nabla v|^2 = \bigg(\frac{\partial v}{\partial x}\bigg)^2 = \bigg(\frac{2\pi}{\varepsilon}\bigg)^2 \cos^2\bigg(\displaystyle\frac{2\pi x}{\varepsilon}\bigg),
\end{equation*}
The numerator is
\begin{equation*}
\begin{split}
\int_{\Omega_\varepsilon}|\nabla v|^2 \,dx\,dy&= \int_{C_\varepsilon}|\nabla v|^2\,dx\,dy = \bigg(\frac{2\pi }{\varepsilon}\bigg)^2\int_{-\frac{\varepsilon^3}{2}}^{\frac{\varepsilon^3}{2}}\int_{-\frac{\varepsilon}{2}}^{\frac{\varepsilon}{2}} \cos^2\bigg(\displaystyle\frac{4\pi x}{\varepsilon}\bigg)\,dx\,dy \\
&= 2\bigg(\frac{2\pi }{\varepsilon}\bigg)^2\varepsilon^3\int_0^{\frac{\varepsilon}{2}}\frac{1+\cos\bigg(\displaystyle\frac{4\pi x}{\varepsilon}\bigg)}{2}\,dx=2\pi^2\varepsilon^2.
\end{split}
\end{equation*}
In this way, since $v$ is zero on $\partial B_{R_1}(x_2)$, we get
\begin{equation*}
	\mu_1(\Omega_\varepsilon)\le \frac{\displaystyle \int_{\Omega_\varepsilon}|\nabla u|^2\,dx}{\displaystyle \int_{\partial \Omega_\varepsilon}u^2 \,d\mathcal{H}^{n-1}}= 2\pi^2\varepsilon,
\end{equation*}
and eventually
\begin{equation*}
	\mu_1(\Omega_\varepsilon)\to 0, \qquad \qquad \text{as}\;\;\varepsilon\to 0. 
\end{equation*}
    \end{proof}

\end{document}
